\subsection{An optional degree-1140 companion}
\label{sec:eleven-optional}

A further positive translation of the same eleven supplied coordinates gives
a second Pareto point. Retain all definitions in Section~\ref{sec:eleven},
put
\[
 \sigma=M(\ell+2c-X)-b,
 \qquad
 W=Lb_0+H\bigl(E(t+\sigma)+k-R-1\bigr),
\]
and change nothing else. The product now has exact degree $1140$ and $129$
binary operations; the shifted two-equation version has residual degrees
$2,1138$ and $124$ operations. Thus the main $123$-operation construction
remains the less expensive one. The optional circuits are separate files.

\paragraph{Soundness of the positive translation.}
The unchanged parity argument gives $N_t=1$ at any product root and hence
all outer bounds. Write $C=C_*$ and $B=\ell+2c-X>0$. Using
$X/q=w\ge1$, direct expansion gives
\begin{align*}
 C\ell-\Delta c
 &=EY\{\eta(X-6-1/q)+\zeta(2X-4)\}-4E\eta\\
 &\hspace{1cm} +(w-1)Y^2c-(4Y+3)c\\
 &\ge EY\{\eta(X-7-1/q)+\zeta(2X-4)\}-(4Y+3)c.
\end{align*}
Since $X,q\ge16$ and $Y\ge4$, both bracket coefficients are at least
$X/2$. As $c<k(Y+1)$ and $(4Y+3)(Y+1)<8Y^2$,
\[
 C\ell-\Delta c>kY^2(X^2/2-8)\ge(k/3)X^2Y^2.
\]
Also $C=XY+2X+2Y+Y/q<3XY/2$. Consequently
\[
 CB-\Delta c>C\ell-\Delta c-CX
 >X^2Y(kY/3-3/2)>1.
\]
It follows that
\[
 \sigma=2TIB+Ri^2c^3(CB-\Delta c)-c-R>0.
\]
Every new positive solution therefore gives a positive old carrier
$t_{\rm old}=t+\sigma$, and the main soundness proof applies directly.

\paragraph{Completeness.}
Subtracting $\sigma$ from the canonical old carrier gives exactly
\[
 t=\gamma+M\{(\tau-\ell)-((a+2)c-g)\}+y+b-h.
\]
The two root differences obey
\[
 \tau-\ell=\frac{\ell k+1}{\tau+\ell}>\frac{k}{3},
 \qquad
 (a+2)c-g=\frac{c^2-1}{(a+2)c+g}
 <\frac{c}{a+2}<\frac{2k}{X+1}<\frac{k}{3}.
\]
Here $\ell>k$, $\tau<\ell+k$, and $c<k(Y+1)$ are enough.
Since $b>h$, the new $t$ is positive and completeness follows.

\paragraph{Exact degree.}
In the predecessor's sixteen-linear-factor expansion the prefactor
$d^8H^{16}$ has degree $464$, and the linear factors have degrees
four times $35$, twice $36$, and ten times $38$. The translation has degree
$35$. For exactly two factors, namely the choices
$g=\sqrt G$, $\tau=\sqrt A$, and $y=-(b-vf)+O(\lambda^{-8})$,
the translated factor becomes
\[
 t+M\{\ell-\sqrt A+(a+2)c-\sqrt G\}-\gamma-vf
    +O(\lambda^{-8}).
\]
Because $\sqrt A-\ell=k/2+O(\lambda^{-9})$, while
$(a+2)c-\sqrt G$ has degree $-1$, its leading term is $-M_0k_0/2$
of degree $28$. The opposite two degree-$35$ factors retain their degree
with twice their original leading term; the other factors are unchanged
in highest degree. The translated predecessor norm therefore has degree
\[
 464+2\cdot28+2\cdot35+2\cdot36+10\cdot38=1042.
\]
The factor $E^{16}$ from index absorption adds $96$; the residual rational
carrier correction has degree $-2$ and affects none of these leaders.
Thus the new norm has degree $1138$ and leading form
\[
 2^{46}i^{48}k_0^{96}w^{90}s_{00}^{188}
 q_0^{700}(q_0-F)^{16}.
\]
The degree-$2$ transport factor yields degree $1140$ for the single equation.
Two complete univariate polynomial checks over different finite fields agree
with this form, and $60$ independent formal conjugate-product checks verify
the optional circuit's norm transcription.
